\section{Заключение}

В ходе лабораторной работы развёрнута и исследована распределённая файловая система GlusterFS на базе четырёх виртуальных машин под управлением Debian~12: одного клиентского узла и трёх серверных. Настроена сеть на основе NAT-сети VirtualBox, каждому узлу присвоен статический IP-адрес. Установлено и запущено ПО GlusterFS~10.3.

Серверные узлы объединены в пул хранения, изучены добавление и удаление узла. Созданы, изменены и удалены тома различных типов: replica, replica с арбитром, distributed. Установлено, что тип stripe удалён из GlusterFS начиная с версии~6, и в качестве его функциональных аналогов исследованы тома disperse (erasure coding) и шардирование. Выполнены добавление и удаление бриков в distributed- и replica-томах, в том числе перебалансировка данных и самолечение реплики.

Настроены квоты на весь том и на отдельный каталог, выявлена задержка их срабатывания и способ её устранения. На LVM-бриках с thin-provisioning создан снапшот тома, выполнено восстановление данных из снапшота и удаление снапшота. Таким образом, все задачи работы выполнены.
